\section{第一轮的观测数与污染敏感性研究}\label{sec:sensitivity}

\subsection{观测数量与污染水平的预登记敏感性分析}
\label{sec:sensitivity-protocol}

为补充单一压力点的评价，另以独立脚本\texttt{sensitivity.py}研究适用范围。
在生成数据前保存的\texttt{research/sensitivity\_plan.json}固定根种子980221，
观测数$n\in\{8,12,16,24\}$、名义污染比例
$\epsilon\in\{0,0.1,0.25,0.4\}$，以及静态异常、恒速异常两个场景。
每个网格单元独立生成12例，合计32格、384例。各格使用由场景编号、观测数和
污染水平派生的独立伪随机子流；三种方法在每格使用完全相同的案例，因此对照差异
按案例配对。固定比较方法仅为最终冻结方案、恒速稳健基线和静态加权最小二乘，
总计1152次预测，不依此重新选择方法或修改参数。

脚本先验证研究冻结、关键源码与开发证据，再生成数据。协议保存预登记文件、
本脚本及冻结方法的摘要；这些本地记录用于审计顺序，不等于第三方时间戳认证。
每格报告平均位置误差、经验CEP90、计算工作量、壁钟时间和实际模型选择频率，
同时保存逐例输入、预测与CV诊断。污染数采用生成器的
$\operatorname{round}(n\epsilon)$，因而小样本下的实际比例具有离散性；
结果必须并列记录名义比例与实际比例。

观测数同时影响信息量与CV门控。默认三折时间交错划分在$n=8$时含训练样本数
为5的折，低于运动模型要求的6份训练观测；若冻结方案采用该CV配置，运动模型会
因可用性条件被排除。该机制可由完成折数与运动可用标记核对，不能将其误读为数据
支持目标静止。若最终方案不含CV，则相关字段记为不适用，不补造门控结果。
本分析仅描述预设合成网格，每格12例不足以精确定位稳健性相变；不作逐格显著性
或多重比较确证，不重新校准误差区域，也不提出新的覆盖率保证。


实际完成384个独立案例、1152次定位预测，耗时0.2分钟。以下结果没有返馈给第一轮研究智能体，也没有用于替换其已冻结方案；第二轮则明确把这些结果作为已暴露先验，另用新数据评价修订。

在全部32个预设网格单元中，所选方法的样本均值低于静态WLS的单元为23个，低于恒速稳健基线的单元为20个。每格仅12例，单元胜出数量是描述性结果，不是显著性检验或部署分布下的成功概率。

\begin{table}[H]\centering\small
\caption{静止目标：所选方法平均位置误差（m）。列为名义污染比例。}

\begin{tabular}{lrrrr}
\toprule
观测数 & 0\% & 10\% & 25\% & 40\%\\
\midrule
8 & 1.807 & 3.042 & 2.810 & 4.937\\
12 & 2.295 & 2.388 & 2.261 & 11.469\\
16 & 1.189 & 2.861 & 2.306 & 3.184\\
24 & 0.897 & 1.506 & 1.940 & 6.133\\
\bottomrule\end{tabular}\end{table}

\begin{figure}[H]\centering
\begin{tikzpicture}\begin{axis}[width=0.91\linewidth,height=6.4cm,xlabel={观测数},ylabel={平均位置误差（m）},xtick={8,12,16,24},ymin=0,grid=major,legend style={at={(0.5,1.04)},anchor=south,legend columns=4}]
\addplot[thick,color=navy,mark=*] coordinates {(8,1.806937) (12,2.295357) (16,1.189295) (24,0.896926)};\addlegendentry{0\%}
\addplot[thick,color=teal,mark=square*] coordinates {(8,3.041615) (12,2.387690) (16,2.861054) (24,1.506256)};\addlegendentry{10\%}
\addplot[thick,color=orange,mark=triangle*] coordinates {(8,2.809855) (12,2.261038) (16,2.305569) (24,1.940370)};\addlegendentry{25\%}
\addplot[thick,color=brick,mark=diamond*] coordinates {(8,4.936753) (12,11.468595) (16,3.184007) (24,6.133056)};\addlegendentry{40\%}
\end{axis}\end{tikzpicture}
\caption{静止目标的样本量与污染敏感性。线段仅连接离散实验格点，不表示格点间的理论插值。}
\label{fig:sensitivity-static_outliers}
\end{figure}

\begin{table}[H]\centering\small
\caption{静止目标：模型选择及相对恒速基线的配对均值差。正差值表示所选方法误差较低。}

\begin{tabular}{lrrrrr}
\toprule
观测数 & 名义污染 & 实际污染 & 选择恒速 & CV可用 & 误差减少/m\\
\midrule
8 & 0\% & 0.0\% & 0.0\% & 0.0\% & 3.582\\
8 & 10\% & 12.5\% & 0.0\% & 0.0\% & 11.203\\
8 & 25\% & 25.0\% & 0.0\% & 0.0\% & 11.262\\
8 & 40\% & 37.5\% & 0.0\% & 0.0\% & 16.244\\
12 & 0\% & 0.0\% & 25.0\% & 100.0\% & 1.646\\
12 & 10\% & 8.3\% & 25.0\% & 100.0\% & 1.097\\
12 & 25\% & 25.0\% & 0.0\% & 100.0\% & 4.197\\
12 & 40\% & 41.7\% & 41.7\% & 100.0\% & 6.950\\
16 & 0\% & 0.0\% & 0.0\% & 100.0\% & 1.332\\
16 & 10\% & 12.5\% & 16.7\% & 100.0\% & 2.038\\
16 & 25\% & 25.0\% & 0.0\% & 100.0\% & 9.356\\
16 & 40\% & 37.5\% & 8.3\% & 100.0\% & 7.023\\
24 & 0\% & 0.0\% & 0.0\% & 100.0\% & 0.964\\
24 & 10\% & 8.3\% & 16.7\% & 100.0\% & 1.237\\
24 & 25\% & 25.0\% & 0.0\% & 100.0\% & 1.827\\
24 & 40\% & 41.7\% & 16.7\% & 100.0\% & 5.043\\
\bottomrule\end{tabular}\end{table}

\begin{table}[H]\centering\small
\caption{移动目标：所选方法平均位置误差（m）。列为名义污染比例。}

\begin{tabular}{lrrrr}
\toprule
观测数 & 0\% & 10\% & 25\% & 40\%\\
\midrule
8 & 11.526 & 9.125 & 10.956 & 14.169\\
12 & 4.647 & 6.090 & 10.801 & 14.661\\
16 & 4.661 & 12.936 & 10.409 & 9.317\\
24 & 1.778 & 2.938 & 5.347 & 8.704\\
\bottomrule\end{tabular}\end{table}

\begin{figure}[H]\centering
\begin{tikzpicture}\begin{axis}[width=0.91\linewidth,height=6.4cm,xlabel={观测数},ylabel={平均位置误差（m）},xtick={8,12,16,24},ymin=0,grid=major,legend style={at={(0.5,1.04)},anchor=south,legend columns=4}]
\addplot[thick,color=navy,mark=*] coordinates {(8,11.526413) (12,4.647328) (16,4.661318) (24,1.777933)};\addlegendentry{0\%}
\addplot[thick,color=teal,mark=square*] coordinates {(8,9.125434) (12,6.090254) (16,12.936288) (24,2.938467)};\addlegendentry{10\%}
\addplot[thick,color=orange,mark=triangle*] coordinates {(8,10.955595) (12,10.801396) (16,10.409263) (24,5.347258)};\addlegendentry{25\%}
\addplot[thick,color=brick,mark=diamond*] coordinates {(8,14.168633) (12,14.661439) (16,9.317081) (24,8.704152)};\addlegendentry{40\%}
\end{axis}\end{tikzpicture}
\caption{移动目标的样本量与污染敏感性。线段仅连接离散实验格点，不表示格点间的理论插值。}
\label{fig:sensitivity-moving_outliers}
\end{figure}

\begin{table}[H]\centering\small
\caption{移动目标：模型选择及相对恒速基线的配对均值差。正差值表示所选方法误差较低。}

\begin{tabular}{lrrrrr}
\toprule
观测数 & 名义污染 & 实际污染 & 选择恒速 & CV可用 & 误差减少/m\\
\midrule
8 & 0\% & 0.0\% & 0.0\% & 0.0\% & -7.328\\
8 & 10\% & 12.5\% & 0.0\% & 0.0\% & 0.496\\
8 & 25\% & 25.0\% & 0.0\% & 0.0\% & 4.017\\
8 & 40\% & 37.5\% & 0.0\% & 0.0\% & 4.096\\
12 & 0\% & 0.0\% & 83.3\% & 100.0\% & -1.647\\
12 & 10\% & 8.3\% & 75.0\% & 100.0\% & -1.341\\
12 & 25\% & 25.0\% & 50.0\% & 100.0\% & -3.113\\
12 & 40\% & 41.7\% & 33.3\% & 100.0\% & 3.313\\
16 & 0\% & 0.0\% & 83.3\% & 100.0\% & -1.208\\
16 & 10\% & 12.5\% & 66.7\% & 100.0\% & -2.377\\
16 & 25\% & 25.0\% & 8.3\% & 100.0\% & -1.925\\
16 & 40\% & 37.5\% & 25.0\% & 100.0\% & -2.142\\
24 & 0\% & 0.0\% & 100.0\% & 100.0\% & -0.000\\
24 & 10\% & 8.3\% & 91.7\% & 100.0\% & -0.240\\
24 & 25\% & 25.0\% & 33.3\% & 100.0\% & -0.098\\
24 & 40\% & 41.7\% & 16.7\% & 100.0\% & -2.198\\
\bottomrule\end{tabular}\end{table}

离散观测数使实际污染比例与名义比例有所差异，表中分别列出。CV可用率为空表示该方法未输出此诊断；不能将缺失信息当作零。若冻结方法使用三折观测门控，则8条观测时部分训练折少于6条，按预定实现会禁止选择恒速；这项设计限制应与由污染或优化失败引起的选择错误区分。

\subsection{门控边界、污染交互与可追溯失败案例}
\label{sec:sensitivity-interpretation}

观测数量改变了信息量，也改变了模型可用性。三折时间交错划分要求每个训练折
至少有6份观测，因此计数条件为$n-\lceil n/3\rceil\geq6$。在$n=8$的96例中，
第一折训练集只有5份观测，程序直接停止CV，完成折数为0，运动可用标记全部为假，
最终全部输出静态模型。该运动无污染单元的均值误差为11.5264 m，恒速稳健
对照为4.1988 m。这里的静态选择来自实现门槛，不能解释成观测证据支持目标静止。

在$n\geq12$的288例中，三折均完成且运动模型全部可用。此时仍存在污染相关的
选择变化：真实运动案例在名义污染率$0,0.1,0.25,0.4$下，选择静态的次数依次为
$4/36,8/36,25/36,27/36$，对应平均位置误差为
3.6955、7.3217、8.8526、10.8942 m。CV比较历史观测上的标准化距离预测损失，
最终任务却评价参考时刻的位置；从距离误差到当前位置误差的放大受轨迹参数、
时间分布和几何条件影响，两类风险没有自动保持大小顺序的保证。

\begin{table}[H]
  \centering
  \small
  \caption{三个实际失败案例的配对位置误差。案例均保留于正式统计，单位为m。}
  \label{tab:sensitivity-failures}
  \begin{tabular}{llrrr}
    \toprule
    案例 & 最终选择 & 所选方法 & 恒速稳健 & 静态WLS\\
    \midrule
    A：运动，$n=16,\epsilon=0.1$ & 恒速 & 64.2834 & 64.2834 & 6.1442\\
    B：运动，$n=16,\epsilon=0.25$ & 静态 & 23.0179 & 3.1658 & 22.8234\\
    C：静态，$n=12,\epsilon=0.4$ & 恒速 & 41.5503 & 41.5503 & 9.3065\\
    \bottomrule
  \end{tabular}
\end{table}

案例A为本次最大误差，求解器仍报告成功、Jacobian秩为4、没有活动数值边界；
最小奇异值约0.16293，条件数约128.22。CV允许并选择恒速，但配对恒速基线也有
相同误差，说明失败不能只归因于模型选择。满秩与优化终止不能保证估计落在真解
附近，也不能代替理论局部误差界的邻域和噪声条件。案例B中运动模型可用，但CV
静态／运动得分为3.20774／3.51454，静态选择比配对恒速结果多出19.8521 m误差。
案例C则把零速度目标解释成运动，直接展示了静态非劣目标仍未得到保证。

为便于复核，表中A、B、C依次对应下列完整\texttt{case\_id}；在
\nolinkurl{results/sensitivity/predictions.json}中按该字段查找即可获得三种方法的
预测和CV诊断，在同目录\texttt{inputs.json}中可核对生成输入：
\begin{description}[leftmargin=1.6em,labelsep=0.4em,font=\normalfont\bfseries]
  \item[A] \nolinkurl{sensitivity-moving_outliers_n16_c1000-r05-921984553727470929}
  \item[B] \nolinkurl{sensitivity-moving_outliers_n16_c2500-r12-13744984418704934917}
  \item[C] \nolinkurl{sensitivity-static_outliers_n12_c4000-r04-14729045951357002381}
\end{description}

本分析为32格、每格12例的探索性描述。不同网格独立抽样，污染数量还经过
$\operatorname{round}(n\epsilon)$取整；观察到的交互和极端案例不能单独识别因果
机制，也不能精确定位污染崩溃阈值。上述解释结合代码规则与保存的诊断提出，
没有据此修改冻结方法、作逐格显著性宣称或建立新的覆盖保证。
